\section{Related Work}

Indirect prompt injection---placing adversarial text where a model-integrated
application will retrieve it, rather than typing it at the model---was
characterized by \citet{greshake2023indirect}, following earlier work on direct
instruction-overriding attacks~\citep{perez2022ignore}. It is catalogued as the
first entry of the OWASP list for generative-AI
applications~\citep{owasp2025llmtop10}.

\ifshort
Benchmarks for tool-integrated agents---InjecAgent~\citep{zhan2024injecagent},
AgentDojo~\citep{debenedetti2024agentdojo}---establish that deployed agents are
susceptible and differ in attack and defense coverage. Our contribution is
orthogonal: how an exfiltration event is \emph{labelled}, and what that label's
false-positive rate is when measured on a matched benign arm rather than
assumed.
\else
Two benchmarks target tool-integrated agents directly.
InjecAgent~\citep{zhan2024injecagent} scores indirect prompt injections against
agents across a large synthetic tool surface, separating attacks that harm the
user from attacks that exfiltrate private data.
AgentDojo~\citep{debenedetti2024agentdojo} provides a dynamic environment in
which attacks and defenses are evaluated against task-specific security
properties rather than a fixed attack list. Both establish that deployed agents
are susceptible; both are concerned with the breadth of attacks and defenses
covered.

Our contribution is orthogonal to that axis. We do not propose new attacks, and
our corpus (\corpusSize{} payloads) is deliberately small. What we propose is a
property of the \emph{label}: that an exfiltration event can be defined so that
determining it requires no judgement, and that the label's false-positive rate
can be measured on a matched adversary-free arm instead of assumed. We are not
aware of a prior agent-injection evaluation that reports a false-positive rate
obtained by running the identical task distribution with every injection
removed; \cref{sec:fp} is our argument that doing so is not optional, because the
natural label fails that test badly.

The planted-artifact idea is itself old: canary values and decoy documents are
long-standing intrusion-detection primitives. The contribution here is not the
canary but its use as a \emph{scoring oracle}---the observation that the
instruction/record asymmetry makes a planted record a sound decision procedure
for agent exfiltration, and that the decision procedure must be scoped to the
destination to be sound in practice.

\ifanon
The agent, workflow, and member-record schema used as the test subject follow a
privacy-constrained contact-center setting from prior benchmarking work on
open-weights versus proprietary models for contact-center tasks, which measures
task capability under a self-hosting constraint and does not consider adversarial
input. The setting is specified in full in \cref{sec:method}, so no part of this
paper depends on that reference; the citation is withheld for anonymity and will
be restored in a non-anonymous version.
\else
The agent, workflow, and member-record schema used as the test subject follow the
privacy-constrained contact-center setting of
\citet{nagaraju2026compliance}, which benchmarks open-weights against
proprietary models on contact-center tasks; that work measures task capability
under a self-hosting constraint and does not consider adversarial input.
\fi
\fi
